% Auto-generated by extract-marking.py -- do not edit by hand unless you know what to do.
% Works for BOTH \documentclass[...,type=solutions,...] and
% [...,type=marking,...] compiles, chosen at compile time via
% \OlympMarkPick / \OlympMarkBeginTable (see olympiad-marking.sty).
% \eqref{n} falls back to a plain '(n)' whenever the referenced
% equation isn't actually typeset in the current document.

\OlympMarkBeginTable{|p{0.88\linewidth}|>{\centering\arraybackslash}p{0.08\linewidth}|}{|p{0.66\linewidth}|>{\centering\arraybackslash}p{0.08\linewidth}|>{\centering\arraybackslash}p{0.06\linewidth}|>{\centering\arraybackslash}p{0.06\linewidth}|}
\hline
\OlympMarkPick{\multicolumn{1}{|c|}{\textbf{Содержание}} & \textbf{Баллы} \\ \hline}{\multicolumn{1}{|c|}{\textbf{Содержание}} & \textbf{Баллы} & \textbf{MP} & \textbf{A} \\ \hline}
\OlympMarkPick{Формула \eqref{1}: $dP = \mu\cdot p dS\cdot \omega r,$ & 0.3 \\ \hline}{Формула \eqref{1}: $dP = \mu\cdot p dS\cdot \omega r,$ & 0.3 & &  \\ \hline}
\OlympMarkPick{\quad\textit{Балл ставится и при записи $\omega r\Leftrightarrow v$ или $dS\Leftrightarrow2\pi rdr\Leftrightarrow rd\varphi dr$ или $pdS\Leftrightarrow dN$. За $P=\mu N v$ или прочие \textbf{не}инфинитезимальные записи балл не ставится.} &  \\ \hline}{\quad\textit{Балл ставится и при записи $\omega r\Leftrightarrow v$ или $dS\Leftrightarrow2\pi rdr\Leftrightarrow rd\varphi dr$ или $pdS\Leftrightarrow dN$. За $P=\mu N v$ или прочие \textbf{не}инфинитезимальные записи балл не ставится.} &  & &  \\ \hline}
\OlympMarkPick{Формула \eqref{2}: $P = \frac23 \mu p\omega \pi R^3.$ & 0.5 \\ \hline}{Формула \eqref{2}: $P = \frac23 \mu p\omega \pi R^3.$ & 0.5 & &  \\ \hline}
\OlympMarkPick{\quad\textit{Вместо коэффициента $2/3$ стоит $1$ либо $1/2$. Если за \eqref{2} стоит \textbf{не}нулевой балл, то ошибка в \eqref{2} должна быть учтена в \eqref{6}, \eqref{7}, и \eqref{10}} & \textit{-0.3} \\ \hline}{\quad\textit{Вместо коэффициента $2/3$ стоит $1$ либо $1/2$. Если за \eqref{2} стоит \textbf{не}нулевой балл, то ошибка в \eqref{2} должна быть учтена в \eqref{6}, \eqref{7}, и \eqref{10}} & \textit{-0.3} & &  \\ \hline}
\OlympMarkPick{Формула \eqref{3}: $j = -\kappa \pi R^2\frac{dT}{dx},$ . Знак минус необязателен, однако это влияет на наличие знака минус в \eqref{5}. & 0.3 \\ \hline}{Формула \eqref{3}: $j = -\kappa \pi R^2\frac{dT}{dx},$ . Знак минус необязателен, однако это влияет на наличие знака минус в \eqref{5}. & 0.3 & &  \\ \hline}
\OlympMarkPick{Формула \eqref{4}: $dq = \alpha\cdot 2\pi R dx\cdot(T-T_0).$ & 0.3 \\ \hline}{Формула \eqref{4}: $dq = \alpha\cdot 2\pi R dx\cdot(T-T_0).$ & 0.3 & &  \\ \hline}
\OlympMarkPick{Формула \eqref{5}: $j(x+dx)-j(x)=dj=-dq.$ . Отсутствие знака минус \textbf{не} принимается, если это приводит к несоответствию в знаках минус между \eqref{3} и этим уравнением. & 0.3 \\ \hline}{Формула \eqref{5}: $j(x+dx)-j(x)=dj=-dq.$ . Отсутствие знака минус \textbf{не} принимается, если это приводит к несоответствию в знаках минус между \eqref{3} и этим уравнением. & 0.3 & &  \\ \hline}
\OlympMarkPick{Формула \eqref{6}: $\frac{d^2T}{dx^2} = \frac{2\alpha}{\kappa R}(T-T_0).$ & 0.2 \\ \hline}{Формула \eqref{6}: $\frac{d^2T}{dx^2} = \frac{2\alpha}{\kappa R}(T-T_0).$ & 0.2 & &  \\ \hline}
\OlympMarkPick{Формула \eqref{7}: $\lambda=\sqrt\frac{2\alpha}{\kappa R},$ & 0.2 \\ \hline}{Формула \eqref{7}: $\lambda=\sqrt\frac{2\alpha}{\kappa R},$ & 0.2 & &  \\ \hline}
\OlympMarkPick{Формула \eqref{8}: $T(x)=T_0+A\exp\ab(-\lambda|x|)$ & 0.4 \\ \hline}{Формула \eqref{8}: $T(x)=T_0+A\exp\ab(-\lambda|x|)$ & 0.4 & &  \\ \hline}
\OlympMarkPick{\quad\textit{Модуль необязателен. Принимается если вместо $T_0$ стоит произвольная \textbf{ненулевая} константа. В остальных случаях балл не ставится.} &  \\ \hline}{\quad\textit{Модуль необязателен. Принимается если вместо $T_0$ стоит произвольная \textbf{ненулевая} константа. В остальных случаях балл не ставится.} &  & &  \\ \hline}
\OlympMarkPick{Формула \eqref{9}: $j(0)=\frac P2.$ & 0.3 \\ \hline}{Формула \eqref{9}: $j(0)=\frac P2.$ & 0.3 & &  \\ \hline}
\OlympMarkPick{Формула \eqref{10}: $\omega = \frac3{\mu pR}\sqrt\frac{2\kappa\alpha }R\cdot(T^*-T_0)=\qty{51.0}{\radian\per\s}.$ & 0.5 \\ \hline}{Формула \eqref{10}: $\omega = \frac3{\mu pR}\sqrt\frac{2\kappa\alpha }R\cdot(T^*-T_0)=\qty{51.0}{\radian\per\s}.$ & 0.5 & &  \\ \hline}
\OlympMarkPick{Численный ответ в \eqref{10} & 0.2 \\ \hline}{Численный ответ в \eqref{10} & 0.2 & &  \\ \hline}
\OlympMarkPick{\quad\textit{Переноса ошибки в численных расчётах нет} &  \\ \hline}{\quad\textit{Переноса ошибки в численных расчётах нет} &  & &  \\ \hline}
\OlympMarkPick{\textbf{Итого} & \textbf{3.5} \\ \hline}{\textbf{Итого} & \textbf{3.5} &  &  \\ \hline}
\end{tabularx}
